\lecture{38}{Нейронные сети и их обучение}

\sourcevideo
    {https://www.youtube.com/playlist?list=PLOzRYVm0a65fhKDS8cYIC7YCuVGJ4FOJx}
    {Week 10: Lecture 38}

Нейронная сеть задаёт параметрическое нелинейное отображение, а её
обучение сводится к минимизации эмпирического риска. При распределении
данных между агентами эта задача снова принимает форму суммы локальных
функций, но размерность параметра и стоимость передачи градиентов
становятся существенными.

\section{Обучение с учителем и один слой}

Пусть задан набор наблюдений
\[
    \mathcal D
    =
    \{(\xi_s,y_s)\}_{s=1}^{M},
    \qquad
    \xi_s\in\mathbb R^{n_0}.
\]
В обучении с учителем требуется построить отображение
\[
    \xi\longmapsto \widehat y=f_\theta(\xi),
\]
где \(\theta\) объединяет все параметры модели. В регрессии выход
обычно принадлежит \(\mathbb R^q\), а в классификации задаёт оценки
вероятностей классов.

Один полносвязный слой преобразует вход \(x\in\mathbb R^n\) по
формуле
\[
    h=\sigma(Wx+b),
\]
где
\[
    W\in\mathbb R^{m\times n},
    \qquad
    b\in\mathbb R^m,
\]
а функция \(\sigma\) применяется покоординатно. Число параметров
слоя равно \(mn+m\).

Смещение можно включить в матрицу весов. Для
\[
    \widetilde x
    =
    \begin{pmatrix}
        x\\
        1
    \end{pmatrix},
    \qquad
    \widetilde W
    =
    \begin{pmatrix}
        W & b
    \end{pmatrix}
\]
имеем
\[
    Wx+b=\widetilde W\widetilde x.
\]
Поэтому термин «веса» иногда используют для всех параметров слоя.

\section{Многослойная сеть и нелинейность}

Полносвязная сеть с \(L\) слоями задаётся рекурсией
\[
    h^0=\xi,
\]
\[
    h^\ell
    =
    \sigma_\ell
    \left(
        W^\ell h^{\ell-1}+b^\ell
    \right),
    \qquad
    \ell=1,\ldots,L-1,
\]
и выходным слоем
\[
    z=W^Lh^{L-1}+b^L.
\]
Векторы \(h^1,\ldots,h^{L-1}\) называются скрытыми
представлениями, а полный набор параметров равен
\[
    \theta
    =
    (W^1,b^1,\ldots,W^L,b^L).
\]

Без функций активации композиция нескольких слоёв остаётся аффинной.
Например,
\[
    W^2(W^1x+b^1)+b^2
    =
    W^2W^1x+W^2b^1+b^2.
\]
Следовательно, увеличение числа линейных слоёв не расширяет класс
представимых функций. Нелинейность вводится именно для устранения
этого вырождения.

Основной пример — функция
\[
    \operatorname{ReLU}(t)=\max\{0,t\}.
\]
Она непрерывна и кусочно-линейна, но не дифференцируема в нуле; при
обучении там выбирают допустимое значение субградиента.

Ширина определяется числом нейронов в слое, а глубина — числом
последовательных преобразований. Глубокая архитектура в некоторых
классах задач может представлять функцию существенно меньшим числом
параметров, чем неглубокая, однако это зависит от класса функций и
требуемой точности. Теоремы универсальной аппроксимации утверждают,
что при подходящей нелинейной активации и достаточном числе нейронов
можно для непрерывной функции \(g\) на компактном множестве \(K\)
добиться
\[
    \sup_{x\in K}
    \lVert f_\theta(x)-g(x)\rVert
    <\varepsilon.
\]
Точная версия результата зависит от архитектуры и выбранной функции
активации.

\section{Выходы регрессии и классификации}

В скалярной регрессии последний слой обычно имеет вид
\[
    \widehat y=w^{\mathsf T}h+b.
\]
Для многомерной регрессии используется несколько выходных координат.

В классификации по \(K\) классам последний линейный слой выдаёт
логиты
\[
    z=(z_1,\ldots,z_K)^{\mathsf T}.
\]
Они не обязаны быть неотрицательными и не суммируются в единицу.
Вероятности получают с помощью softmax:
\[
    p_r
    =
    \frac{\exp(z_r)}
    {\displaystyle\sum_{q=1}^{K}\exp(z_q)},
    \qquad
    r=1,\ldots,K.
\]
Тогда
\[
    p_r>0,
    \qquad
    \sum_{r=1}^{K}p_r=1,
\]
а предсказываемый класс равен
\[
    \widehat r
    =
    \operatorname*{arg\,max}_{r}p_r.
\]

Для численно устойчивого вычисления положим
\[
    m=\max_r z_r.
\]
Поскольку общий множитель \(\exp(-m)\) сокращается,
\[
    p_r
    =
    \frac{\exp(z_r-m)}
    {\displaystyle\sum_{q=1}^{K}\exp(z_q-m)}.
\]
Теперь \(z_r-m\le0\), поэтому наибольшая экспонента равна единице и
переполнение предотвращается.

\section{Эмпирический риск и функции потерь}

Потеря на одном объекте записывается как
\[
    \ell_s(\theta)
    =
    \ell\bigl(f_\theta(\xi_s),y_s\bigr).
\]
Эмпирический риск равен
\[
    \mathcal L(\theta)
    =
    \frac1M\sum_{s=1}^{M}\ell_s(\theta),
\]
а обучение формулируется задачей
\[
    \min_\theta \mathcal L(\theta).
\]

Для регрессии часто используется среднеквадратичная потеря
\[
    \ell_s(\theta)
    =
    \frac12
    \lVert f_\theta(\xi_s)-y_s\rVert^2,
\]
откуда
\[
    \mathcal L_{\mathrm{MSE}}(\theta)
    =
    \frac1{2M}
    \sum_{s=1}^{M}
    \lVert f_\theta(\xi_s)-y_s\rVert^2.
\]

В классификации правильный ответ задаётся индикаторным вектором
\[
    t=(t_1,\ldots,t_K)^{\mathsf T},
\]
а потеря определяется кросс-энтропией
\[
    \ell_{\mathrm{CE}}
    =
    -\sum_{r=1}^{K}t_r\log p_r.
\]
Если правильный класс имеет номер \(r^\star\), то
\[
    \ell_{\mathrm{CE}}=-\log p_{r^\star}.
\]
Минимизация этой величины увеличивает вероятность правильного класса.

Квадратичная потеря выпукла по непосредственному предсказанию, а
кросс-энтропия с softmax имеет хорошие свойства по логитам. Однако
отображение \(\theta\mapsto f_\theta(\xi)\) в многослойной сети
нелинейно и содержит композиции параметров. Поэтому
\(\mathcal L(\theta)\) в общем случае невыпукла, и градиентные методы
не гарантируют нахождение глобального минимума произвольной сети.

\section{Режимы вычисления градиента}

Полный градиент эмпирического риска равен
\[
    \nabla\mathcal L(\theta)
    =
    \frac1M
    \sum_{s=1}^{M}\nabla_\theta\ell_s(\theta),
\]
а шаг полного градиентного спуска имеет вид
\[
    \theta^{k+1}
    =
    \theta^k
    -
    \eta_k
    \frac1M
    \sum_{s=1}^{M}
    \nabla_\theta\ell_s(\theta^k).
\]
Одна итерация использует весь набор данных и потому даёт стабильное,
но дорогое направление.

В SGD выбирается случайный индекс \(s_k\):
\[
    \theta^{k+1}
    =
    \theta^k
    -
    \eta_k\nabla_\theta\ell_{s_k}(\theta^k).
\]
При равномерной выборке оценка несмещённая:
\[
    \mathbb E
    \left[
        \nabla_\theta\ell_{s_k}(\theta)
    \right]
    =
    \nabla\mathcal L(\theta).
\]
Итерация дешева, но дисперсия направления может быть большой.

Для мини-пакета \(B_k\subset\{1,\ldots,M\}\)
\[
    g_k
    =
    \frac1{|B_k|}
    \sum_{s\in B_k}
    \nabla_\theta\ell_s(\theta^k),
    \qquad
    \theta^{k+1}=\theta^k-\eta_kg_k.
\]
Увеличение \(|B_k|\) обычно уменьшает дисперсию, но повышает стоимость
итерации. Полный градиент соответствует \(|B_k|=M\), а чистый SGD —
\(|B_k|=1\). Мини-пакеты дают компромисс и позволяют эффективно
использовать параллельные вычислители.

Шум SGD иногда помогает не задерживаться в неглубоких локальных
областях, но это не является общей гарантией нахождения лучшего
минимума. Поведение зависит от геометрии функции, шага, размера
мини-пакета и структуры шума; слишком большая дисперсия вызывает
колебания и замедляет сходимость.

\section{Распределённое обучение}

Пусть данные разделены между агентами:
\[
    \mathcal D
    =
    \bigcup_{i=1}^{N}\mathcal D_i.
\]
Локальный риск имеет вид
\[
    f_i(\theta)
    =
    \frac1{|\mathcal D_i|}
    \sum_{(\xi,y)\in\mathcal D_i}
    \ell\bigl(f_\theta(\xi),y\bigr),
\]
а глобальная задача записывается как
\[
    \min_\theta
    \sum_{i=1}^{N}\omega_i f_i(\theta),
\]
где веса \(\omega_i\) могут учитывать размеры локальных наборов. Это
та же структура суммы локальных функций, что и в распределённой
оптимизации, но градиенты обычно вычисляются стохастически, а
размерность параметра может быть огромной.

Пусть число параметров равно \(d_\theta\). Передача полного вектора
параметров или градиента требует сообщения размерности \(d_\theta\).
При обмене центрального сервера с \(N\) агентами объём одного раунда
имеет порядок
\[
    O(Nd_\theta).
\]
Поэтому при миллионах или миллиардах параметров коммуникация может
оказаться дороже локального вычисления. Архитектура распределённого
обучения должна учитывать расположение данных, способ агрегации,
синхронность, топологию сети и объём сообщений. Эти вопросы являются
предметом следующих лекций.
